% Nová příprava podle aktuálního GitHubu Work-MAT, hodina 93 (nahrazuje vyjadrovani ze vzorce).
% Realne vztahy vs modelova data jsou odlisena; chemicke reakce jsou jen teoreticke vycisleni.
\documentclass[11pt,a4paper]{article}
\input{../../sablony/pripravy/priprava.sty}
\nastavPripravu{1. ročník}{Rovnice a nerovnice}{93}{Aplikační slovní úlohy napříč obory}
\begin{document}
\blok{Motivace: matematika z reálných oborů}
Žáci už umějí upravovat vzorce. Dnes je \textbf{nebudeme znovu vyjadřovat izolovaně}: v různých situacích si vybereme matematický nástroj a vysvětlíme, proč je přiměřený. Pracujeme s jednotkami, omezeními modelu a s věcnou kontrolou.
\poznamka{Tato příprava je výběrová, nikoli sedm povinných příkladů do 45 minut. Doporučený plán: 5 min úvod, 2--3 modelové problémy po 10--12 min, 8--10 min samostatná práce a reflexe. Ostatní příklady jsou hotová rezerva.}
\blok{Jak poznat vhodnou matematickou metodu}
\begin{itemize}
\item Jedna neznámá a rovnost: \textbf{lineární rovnice}; požadovaná mez: \textbf{nerovnice}.
\item Dvě neznámé a dvě nezávislé podmínky: \textbf{soustava rovnic}.
\item Vzdálenost, odchylka, tolerance: \textbf{absolutní hodnota}.
\item Hodnota ovlivňující počet kořenů: \textbf{parametr}.
\item Zachování počtu atomů v reakci: \textbf{lineární bilance / soustava} (volitelné chemické rozšíření).
\end{itemize}
\poznamka{Čísla v úlohách (s výjimkou citovaných fyzikálních a chemických konstant a skutečných vztahů) jsou didaktické modely. Zejména zpoždění NAND/XOR se mění podle technologie, zátěže a zapojení; z jednoho příkladu nelze vyvozovat univerzální žebříček rychlosti hradel.}
\blok{Výběr sedmi aplikačních příkladů (jeden na typ)}
\begin{prikladReseniBox}[title={Lineární rovnice -- ELE}]
Zdroj má napětí $12\,\mathrm V$. LED v pracovním bodě vykazuje úbytek $2\,\mathrm V$ a v sérii je rezistor $500\,\Omega$. Podle modelu $2+500I=12$ určete proud $I$ v ampérech a miliampérech.
\tcblower
\[\begin{aligned}
2+500I&=12\\500I&=10\\I&=0{,}020\,\mathrm A=20\,\mathrm{mA}
\end{aligned}\]
\par\textit{Poznámka:} Jde o model s přibližně konstantním úbytkem napětí na LED.
\end{prikladReseniBox}
\begin{prikladReseniBox}[title={Lineární nerovnice -- MEC}]
Plošina má modelovou nosnost $360\,\mathrm{kg}$. Nářadí váží $60\,\mathrm{kg}$ a každý pracovník $75\,\mathrm{kg}$. Kolik pracovníků se může bezpečně vejít podle hmotnostního limitu? Neřešte další bezpečnostní předpisy.
\tcblower
\[\begin{aligned}
60+75n&\le360\\75n&\le300\\n&\le4\\n_{\max}&=4\ \text{osoby}
\end{aligned}\]
\end{prikladReseniBox}
\begin{prikladReseniBox}[title={Soustavy lineárních rovnic -- ICT}]
Ve stejné modelové technologii byly změřeny dvě sériové cesty: dva NAND a jeden XOR mají dohromady $14\,\mathrm{ns}$, jeden NAND a dva XOR $16\,\mathrm{ns}$. Určete z těchto měření zpoždění $n$ jednoho NAND a $x$ jednoho XOR a porovnejte je.
\tcblower
\[\begin{aligned}
2n+x&=14\\n+2x&=16\\x&=14-2n\\n+2(14-2n)&=16\\n&=4\,\mathrm{ns},\quad x=6\,\mathrm{ns}\\K&=\{[4;6]\}
\end{aligned}\]
\par\textit{Poznámka:} Výsledek platí pouze pro tuto modelovou sadu měření, zpoždění závisí na technologii a zátěži.
\end{prikladReseniBox}
\begin{prikladReseniBox}[title={Rovnice s absolutní hodnotou -- MEC}]
Robot je na přímé kolejnici v poloze $x$ (metry). Kontrolní bod je na souřadnici $50\,\mathrm m$. Robot má být přesně $8\,\mathrm m$ od něj. Určete všechna $x$ z $|x-50|=8$.
\tcblower
\[\begin{aligned}
|x-50|&=8\\x-50&=8\ \text{nebo}\ x-50=-8\\x&=58\ \text{nebo}\ x=42\\K&=\{42;58\}\,\mathrm m
\end{aligned}\]
\end{prikladReseniBox}
\begin{prikladReseniBox}[title={Nerovnice s absolutní hodnotou -- ELE}]
Zdroj má jmenovité napětí $12\,\mathrm V$ a povolenou odchylku nejvýše $0{,}5\,\mathrm V$. Určete interval přijatelných hodnot $U$ z $|U-12|\le0{,}5$.
\tcblower
\[\begin{aligned}
|U-12|&\le0{,}5\\-0{,}5\le U-12&\le0{,}5\\11{,}5\le U&\le12{,}5\\K&=\langle11{,}5;12{,}5\rangle\,\mathrm V
\end{aligned}\]
\end{prikladReseniBox}
\begin{prikladReseniBox}[title={Lineární rovnice s parametrem -- ELE}]
V idealizované elektrické větvi platí $(a-2)I=6$, kde $a$ je reálný parametr odpovídající odporu v ohmech a $I$ je orientovaný proud v ampérech. Diskutujte řešení pro všechna $a$.
\tcblower
\[\begin{aligned}
(a-2)I&=6\\a\ne2:&\ I=\frac6{a-2}\\a=2:&\ 0=6\ \Rightarrow\ K_2=\varnothing\\K_a&=\begin{cases}\{\frac6{a-2}\},&a\ne2,\\\varnothing,&a=2.\end{cases}
\end{aligned}\]
\par\textit{Fyzikální interpretace:} U běžného pasivního odporu požadujeme $a-2>0$, tedy $a>2$. Záporné matematické hodnoty odporu nelze automaticky vykládat jako běžný rezistor.
\end{prikladReseniBox}
\begin{prikladReseniBox}[title={Chemické rovnice jako soustavy bilancí -- CHE}]
Doplňte nejmenší kladné celočíselné koeficienty $a,b,c$ do chemické rovnice $a\,\mathrm{CaO}+b\,\mathrm{H_2O}\rightarrow c\,\mathrm{Ca(OH)_2}$. Sestavte bilanci atomů Ca, O a H; vysvětlete, proč lze jeden koeficient zvolit.
\tcblower
\[\begin{aligned}
\text{Ca: }a&=c\\\text{O: }a+b&=2c\\\text{H: }2b&=2c\\a&=b=c=1\\\mathrm{CaO+H_2O}&\rightarrow\mathrm{Ca(OH)_2}
\end{aligned}\]
\par\textit{Poznámka:} Reakce je skutečná a exotermická; nejde o návod na praktické provádění.
\end{prikladReseniBox}
\blok{Samostatná práce -- volba podle zaměření třídy}
\begin{zadaniReseniBox}[title={Vyber dva úkoly z různých oborů}]
Otevři přiloženou \textbf{sbírku 93} a vyber dva různé typy matematiky. Doporučené dvojice: \textbf{SOU-03 + ABN-02} (elektrotechnika a mechanika), \textbf{ROV-03 + NER-03} (ICT), nebo \textbf{ABS-05 + PAR-05} (doprava). Nejdřív sestav vlastní model, pak zkontroluj omezení a jednotky.
\tcblower
Postupy, mezivýsledky i odpovědi jsou v samostatném souboru \texttt{93\_SB\_AplikacniUlohy\_reseni.pdf}. Učitel vybere úlohy podle oboru konkrétní třídy. K hodnocení nestačí samotné číslo; žák musí popsat předpoklady modelu.
\end{zadaniReseniBox}
\blok{Domácí úkol}
\begin{zadaniReseniBox}[title={Dobrovolné rozšíření -- sbírka 49 úloh}]
Vyber si \textbf{alespoň dvě dobrovolné úlohy} z různých matematických kategorií a rozdílných oborů. Není nutné vyřešit vše; volitelně můžeš zkusit více oddílů. U každé uveď matematický model, výsledek s jednotkami a zkoušku.
\tcblower
Využij učitelský klíč. Cílem je argumentace, nikoli mechanické dosazování. Zvlášť sleduj, které modelové předpoklady by se v praxi musely ověřit.
\end{zadaniReseniBox}
\blok{Závěr a reflexe}
\poznamka{Závěrečná otázka: Kde přestává model odpovídat realitě? Příklad: GHz není totéž co miliardy hotových instrukcí za sekundu; u směsí předpokládáme zachování hmotnosti; u počtu lidí musí vyjít celé nezáporné číslo. Chemické vyčíslování neprovádějte laboratorně.}
\end{document}